\documentclass[10pt,conference]{IEEEtran}
\IEEEoverridecommandlockouts
% The preceding line is only needed to identify funding in the first footnote. If that is unneeded, please comment it out.
\usepackage{cite}
\usepackage{amsmath,amssymb,amsfonts}
\usepackage{algorithmic}
\usepackage{graphicx}
\usepackage{textcomp}
\usepackage{xcolor}
\usepackage{booktabs}
\usepackage{url}
\usepackage{tikz}
\usetikzlibrary{arrows.meta,positioning,shapes.geometric,fit,backgrounds}

\def\BibTeX{{\rm B\kern-.05em{\sc i\kern-.025em b}\kern-.08em
    T\kern-.1667em\lower.7ex\hbox{E}\kern-.125emX}}

% ── Provisional numeric values ───────────────────────────────────────────────
% Wrap every tweakable statistic in \dnum{...}. While \shownumstrue, these
% render highlighted (colored) so they are easy to find and update as the data
% grows. To remove all highlighting for a clean PDF, flip to \shownumsfalse.
\newif\ifshownums\shownumstrue
\newcommand{\dnum}[1]{\ifshownums\textcolor{red!75!black}{\textbf{#1}}\else#1\fi}

\newcommand{\researchquestion}[2]{%
  \vspace{0.4em}
  \noindent\colorbox{gray!15}{%
    \parbox{\dimexpr\linewidth-2\fboxsep}{%
      \vspace{0.3em}%
      \textbf{#1:} #2%
      \vspace{0.3em}%
    }%
  }%
  \vspace{0.4em}%
}

\begin{document}

\title{SecDevQA: Can LLM Agents Answer Security Questions Developers Ask?}

\author{
\IEEEauthorblockN{Anonymous Authors}
}
\maketitle

\begin{abstract}
Developers routinely ask security questions in open-source project forums: whether a disclosed vulnerability affects their specific usage, which release fixes it, or whether a proposed patch is sufficient. Answering such questions requires cross-referencing artifacts that live partly outside the codebase---advisories, fix commits, dependency manifests---yet no existing benchmark evaluates LLMs or coding agents on the security questions developers actually ask, nor on the artifact dependence those questions exhibit. We present \mbox{SecDevQA}, a benchmark of \dnum{[N]} developer security question--answer pairs mined from issue threads across \dnum{[K]} open-source repositories and restricted to pairs whose maintainer answer states at least one externally verifiable fact (CVE/GHSA identifier, fixed version, fix commit or pull request, advisory record). Each verified thread is normalized into a self-contained, resolution-free question and a thread-grounded gold answer, labelled by \emph{knowledge type}---whether a correct answer is derivable from general security knowledge (parametric) or requires project-specific or advisory-specific information (grounded)---and annotated with the artifact types the maintainer drew on. Systems are evaluated at the moment the question was asked: an agent receives tool access over a \emph{time-capped snapshot} of the project (repository at the last pre-report commit; issues, pull requests, and advisories published before the report), so that answering is reconstruction of the maintainer's reasoning rather than retrieval of the posted resolution. Grading is two-tiered and condition-aware: a deterministic identifier check scores each verifiable fact only under conditions where it is knowable, and a per-claim LLM judge---validated against human labels---scores the free-text remainder, with fabricated specifics reported separately as hallucinations. Analysis of the corpus yields a taxonomy of developer security question categories and an empirical map from categories to required artifact types; selectively granting and withholding artifact access to frontier LLMs, open-weight LLMs, and coding agents quantifies the marginal value of each artifact type per category. SecDevQA, the mining pipeline, and the evaluation harness are released as open artifacts.
\end{abstract}

\begin{IEEEkeywords}
software security, developer questions, benchmark, large language models, coding agents, question answering, mining software repositories
\end{IEEEkeywords}

\section{Introduction}
\label{sec:intro}

Developers regularly confront security decisions in the course of building and maintaining their own software. They need to know whether a freshly disclosed vulnerability in a dependency actually affects the way they use it, which release first introduced a regression, whether a proposed patch is sufficient, or which version is safe to upgrade to. When the answer is not obvious, they ask---on issue trackers, in advisory discussions, and on CVE-linked threads of the projects they depend on. How developers resolve such questions has direct consequences for the security of the software they ship: a substantial body of usable-security research shows that developers lean heavily on online information sources when solving security problems, and that the quality of those sources measurably shapes the security of the resulting code. Acar et al.\ found that only a small fraction of the Stack~Overflow threads developers consult for security tasks contain secure advice~\cite{acar2016information}, and Fischer et al.\ traced insecure snippets from such threads into hundreds of thousands of shipped applications~\cite{fischer2017stackoverflow}. Getting trustworthy answers to developer security questions is therefore not a convenience problem but a software-security problem.

These questions are also distinctive in what it takes to answer them. Unlike general code-comprehension questions, a developer security question typically cannot be resolved from the repository's source alone. Answering ``does CVE-2022-25883 affect me, and what do I upgrade to?'' requires cross-referencing artifacts that live partly outside the codebase: the CVE or GHSA advisory and its affected-version ranges, the fix commit or pull request, the project's dependency manifest, and sometimes an upstream RFC or registry. Project maintainers answer these questions precisely because they know which artifact to consult---and, crucially, \emph{different questions demand different artifacts}. A question about a transitive dependency is answered from the advisory and the manifest; a question about a TLS validation change is answered from the source and the commit history. Which artifact types are necessary for which kinds of security question is, to date, an unmeasured empirical question.

Large language models (LLMs) and coding agents are increasingly the first responders to such questions, yet no existing benchmark evaluates them on the security questions developers actually ask, or on the artifact dependence those questions exhibit. Repository-level QA benchmarks---SWE-QA~\cite{sweqa2026}, CoReQA~\cite{coreqa2025}, and StackRepoQA~\cite{stackrepoqa2025}---treat developer questions uniformly and consider only code-internal context; they neither isolate security questions nor model the external advisory artifacts those questions depend on. (StackRepoQA further reports that much of the apparent accuracy on such tasks reflects memorization of public threads rather than reasoning~\cite{stackrepoqa2025}, a contamination risk any security-QA benchmark must confront.) Security-specific LLM benchmarks, in turn, measure a different task: SecVulEval~\cite{secvuleval2025} and CodeSecEval~\cite{codeseceval2024} evaluate vulnerability detection and secure code generation, SecureAgentBench~\cite{secureagentbench2025} evaluates whether agents write secure patches, and SecQA~\cite{secqa2023} tests textbook security knowledge via multiple-choice questions. None of them studies real, free-text developer security questions, and none treats artifact context as a controlled variable.

We close this gap with \textbf{SecDevQA}, a benchmark of real developer security question--answer pairs mined from GitHub issues, discussions, and CVE-linked threads across open-source GitHub repositories. Two design choices make the benchmark both rigorous and analytically useful. First, we restrict the benchmark to \emph{hard-verifiable} pairs: pairs whose maintainer answer states at least one externally checkable fact---a CVE or GHSA identifier, a fixed version number, a fix commit or pull request, or an advisory URL. These hard facts anchor automated grading, turning the otherwise subjective task of scoring a free-text answer into the narrowly scoped task of checking whether a model's response correctly conveys a named, verifiable claim---a judgment a human can validate in seconds. Second, every pair is annotated with the artifact types its answer drew on, so that artifact context can be supplied to models as a controlled experimental condition rather than left implicit.

This design lets us pursue three questions that prior work cannot. We first derive, by inductive open coding of the mined corpus, a taxonomy of developer security question categories and an empirical map of which artifact types maintainers draw on to answer each---findings about developer security practice that stand independently of any model evaluation (RQ1). We then measure how accurately frontier LLMs, open-weight LLMs, and coding agents answer these questions, broken down by category (RQ2). Finally, and centrally, we ask which artifact types are \emph{necessary} to answer which categories of question, by supplying artifacts as context under controlled conditions and measuring each type's marginal contribution (RQ3); we further examine whether autonomous agents
retrieve the artifacts they actually need (RQ4). Our guiding hypothesis is that context dependence is category-specific: an artifact type that is essential for one category of security question yields little benefit for another---a result with direct implications for how retrieval-augmented and agentic security tools should be built.

\noindent\textbf{Contributions.} This paper makes the following contributions:
\begin{itemize}
  \item \textbf{SecDevQA}, a benchmark of real, hard-verifiable developer security Q\&A pairs mined from open-source repositories, each annotated with the artifact types needed to answer it and a structured record of its verifiable hard facts (\S\ref{sec:benchmark-construction}).
  \item A fact-grounded grading protocol that exploits hard facts to make LLM-judge scoring narrowly scoped and human-validatable
        (\S\ref{sec:benchmark-construction}).
  \item An inductively derived taxonomy of developer security question categories and an empirical category-to-artifact map, both independent of any LLM evaluation (\S\ref{sec:open-coding}).
  \item An evaluation of frontier and open-weight LLMs and coding agents under controlled context conditions, quantifying the marginal, category-specific contribution of each artifact type (\S\ref{sec:experimental-setup}).
  \item A public release of the benchmark, the mining and extraction pipeline, and the evaluation harness.
\end{itemize}

\section{Benchmark Construction}
\dnum{[The numeric data throughout the paper will be iteratively updated, hence they are colored and bold]}
\label{sec:benchmark-construction}

SecDevQA is built in five stages: (i)~systematic selection of source
repositories from the GitHub Advisory Database, (ii)~LLM-assisted detection and
extraction of candidate security question--answer (Q\&A) pairs from issue and
discussion threads, (iii)~manual verification and artifact/hard-fact annotation
of each candidate pair, (iv)~normalization of every verified pair into one or
more self-contained, gradeable evaluation items, and (v)~inductive open coding of
the verified pairs toward a taxonomy of developer security questions. The numbers reported below describe the corpus at the
current stage of construction; the benchmark is being actively expanded along
the same protocol, so absolute counts will grow while the procedure remains
fixed.

\subsection{Repository Selection}
\label{sec:repo-selection}

To avoid hand-picking projects in a way that could bias the question distribution, we draw the source repositories from an external, independent maintained catalogue of security activity: the GitHub Advisory Database. We clone the database (snapshot \dnum{2026-06-04}; \dnum{31{,}213} reviewed advisories) and parse every reviewed GHSA entry, extracting the GitHub repositories referenced by each advisory and, for each repository, the number of advisories that carry a documented fix version. This yields \dnum{9{,}102} distinct repositories that have appeared in at least one reviewed advisory. We then apply three pre-specified filters intended to keep repositories that are
both security-active and practically minable:

\begin{enumerate}
  \item \textbf{Security activity:} at least \dnum{10} advisories with a
        documented fix version, so that the project has a substantive, externally
        catalogued security history;
  \item \textbf{Popularity:} at least \dnum{1{,}000} GitHub stars, as a proxy for
        a real user base whose developers ask questions;
  \item \textbf{Maintenance:} a push within one week (as of \dnum{06-06-2026}) and not archived, so that threads reflect current practice.
\end{enumerate}

Applying the advisory-count filter leaves \dnum{440} repositories; adding the
popularity and maintenance filters leaves \dnum{344}. To enforce language and
ecosystem diversity rather than concentrating on whichever ecosystem happens to
have the most advisories, we rank the survivors by fix-bearing advisory count
within each package ecosystem (npm, PyPI, RubyGems, etc.) and retain the top
repositories per ecosystem, producing a diverse candidate pool spanning
\dnum{ten} ecosystems.

From this pool we are mining repositories incrementally. The current corpus
comprises the \dnum{ten} repositories in Table~\ref{tab:repos}, spanning
\dnum{three} languages (Python, JavaScript, Ruby) and a range of
security-sensitive domains:
HTTP clients (\texttt{requests}, \texttt{urllib3}, \texttt{axios}),
cryptography (\texttt{pyca/cryptography}), image processing (\texttt{Pillow}),
web frameworks (\texttt{express}, \texttt{fastapi}, \texttt{rails}),
authentication (\texttt{node-jsonwebtoken}), and payment integration
(\texttt{stripe-node}). The selection deliberately mixes repositories with a
large catalogued advisory history with widely used libraries in domains where
security questions are common but formal advisories are sparser, so that the
corpus is not skewed toward projects with heavyweight advisory processes
(\S\ref{sec:verification} describes the per-repository cap that enforces this).

\subsection{Detection and Extraction of Q\&A Pairs}
\label{sec:extraction}

For each repository we retrieve the most recent \dnum{2{,}000} issue and
discussion threads and flatten them into ordered comment sequences. We bound the
window deliberately: our goal is not to harvest every security Q\&A pair a
project has ever produced, but to collect a corpus large and diverse enough to
support the taxonomy and evaluation in this study. Recent threads also better
reflect current libraries, advisories, and developer practice.

Locating security Q\&A pairs in this material by hand is impractical. A single
active repository accumulates thousands of issues, the large majority of which
are ordinary functional bugs, feature requests, or support questions with no
security content; security discussions are a sparse minority scattered across
that volume, and confirming whether a given thread contains a genuine
question--answer exchange requires reading it end to end. Manually browsing
every issue across \dnum{ten} repositories---tens of thousands of threads---to
find this minority is prohibitively slow. We therefore use an LLM as a
\emph{prefilter} that reads every thread and surfaces likely candidates, turning
an intractable full-text triage into a tractable verification problem over a
small candidate set. A two-stage LLM pipeline identifies candidate Q\&A pairs.

\textbf{Stage~1 (detection).} An LLM reads each thread and decides whether it
contains a genuine developer security question paired with a substantive answer,
emitting a boolean, a short summary, and a confidence score. Because this stage
is a prefilter feeding human verification, we deliberately prompt it to favor
\emph{recall} over precision: the cost of a false positive is a few seconds of
reviewer time at the verification stage, whereas a false negative silently drops
a real pair from the corpus with no opportunity for recovery. We accordingly set
a permissive \dnum{0.70} confidence threshold and instruct the model to surface
borderline threads rather than discard them. The downstream \dnum{39\%} rejection
rate (\S\ref{sec:verification}) is the expected, acceptable price of this
recall-oriented design.

\textbf{Stage~2 (extraction).} For threads passing Stage~1, a second prompt
extracts the specific question and answer comments, the role of the answerer
(maintainer, contributor, commenter, or self-answer), the artifact types the
answer draws on (\S\ref{sec:annotation}), and any externally verifiable
\emph{hard facts} the answer states---CVE/GHSA/OSV identifiers, fixed version
numbers, fix commit SHAs, fix PR references, and advisory URLs.

This pipeline produced \dnum{XXXXX} candidate pairs across the \dnum{ten}
repositories. We intentionally extract candidates broadly at this stage and rely
on manual verification, rather than the LLM, to make the final inclusion
decision.

\subsection{Manual Verification}
\label{sec:verification}

Every candidate pair is reviewed by the first author through a purpose-built
annotation interface that presents the full thread alongside the extracted
question, answer, role, artifacts, and hard facts. A pair is \emph{accepted}
only if it satisfies the inclusion criteria: the question concerns a security
situation in the developer's own project, the answer is substantive and
information-providing, and the question and answer are attributable to distinct,
identifiable comments. Pairs that are not genuinely security-related
(e.g.\ ordinary functional bugs), that are unanswerable without information
absent from the thread, or that are mis-extracted are \emph{rejected}; the
reviewer records a short reason for each rejection.

To prevent the most security-active projects from dominating the corpus, we cap
the number of verified pairs retained per repository at \dnum{25}. This keeps the
corpus balanced across projects, languages, and security domains rather than
skewed toward the few repositories with the densest advisory history. In the
current corpus the cap does not yet bind---every repository falls at or below it
(the largest, \texttt{requests}, contributes \dnum{24} pairs)---but it fixes the
composition policy as mining scales to more and larger repositories.

Of the \dnum{XXXXX} candidates, \dnum{158} (\dnum{60.8\%}) were accepted and
\dnum{102} were rejected---a rejection rate that, given the recall-oriented
prefilter, is expected and underscores why LLM extraction alone is insufficient
and human verification is required. No single repository exceeds \dnum{15.2\%} of
the verified pairs. The accepted pairs are answered predominantly by project
maintainers (\dnum{104} of \dnum{158}), with the remainder from
contributors~(\dnum{26}), other commenters~(\dnum{18}), and the original poster
self-answering~(\dnum{10}). Threads span \dnum{2020--2026} and contain a median
of \dnum{4} comments (mean \dnum{5.7}, max \dnum{41}).

Of the \dnum{158} verified pairs, \dnum{105} (\dnum{66.5\%}) satisfy the
additional \emph{hard-verifiability} requirement: the answer states at least one
externally checkable fact. This hard subset is the part of the benchmark that
supports fact-grounded automated grading (the remaining \dnum{53} pairs are
retained for taxonomy construction so that the taxonomy is not biased toward only
those question types that yield checkable identifiers). Across the hard subset,
the most common fact types are fixed version numbers (\dnum{47} pairs) and
advisory URLs (\dnum{44}), followed by fix PRs (\dnum{26}), CVE identifiers
(\dnum{24}), fix commits (\dnum{21}), and GHSA identifiers (\dnum{14}). Hard
facts are spot-checked against NVD, GHSA, and OSV at construction time, and the
advisory-database snapshot date is recorded with the release to support future
re-verification.

\subsection{Artifact-Type and Hard-Fact Annotation}
\label{sec:annotation}

Each verified pair carries an \texttt{artifacts\_needed} list recording which
artifact types the maintainer's answer drew on, and a structured
\texttt{hard\_facts} record. The artifact-type distribution
(Table~\ref{tab:artifacts}) shows that answering these questions is rarely a
matter of reading source code alone: documentation and project source dominate,
but advisories, pull-request and commit history, dependency manifests, and
external references (RFCs, registries, upstream issues) each appear in a
non-trivial fraction of answers. This spread is what motivates the per-artifact
context conditions in the evaluation (\S\ref{sec:experimental-setup}): different
question categories lean on different artifact types, and the annotation makes
that dependency measurable.

\subsection{Question Normalization}
\label{sec:normalization}

A verified thread is not yet a gradeable evaluation item. The developer's
question is embedded in a multi-comment exchange, leans on earlier comments
(``the error above,'' ``this issue''), and frequently sits in a thread whose
later comments---or even whose title---already disclose the resolution. Handing
such a raw thread to a model under test would leak the answer and conflate reading
comprehension with security reasoning. We therefore normalize each verified pair
into one or more self-contained \emph{(question, answer)} items through a single
LLM pass that is then human-verified. Every decision in this step is governed by
one constraint---the item must stay faithful to what was really asked and
answered---because authenticity is the benchmark's reason for being; the choices
below are each justified against that constraint.

\noindent\textbf{Self-contained, resolution-free questions.}
The question is reconstructed to stand alone: every reporter-supplied specific
(versions, runtime and operating system, configuration, the exact error and
stack-trace text, and all code or proof-of-concept, kept verbatim) is preserved,
and dangling references are resolved against earlier comments. The question must
not, however, state its own resolution---the fixed version, patch commit or pull
request, or advisory identifier introduced \emph{as the answer} is withheld---so
that the item remains a genuine open question rather than a fill-in-the-blank. We
separate the \emph{subject} a reporter cites as their premise (the CVE they are
worried about, the version they run), which legitimately stays, from the
\emph{fix} that resolves it, which must not appear. To make this auditable rather
than a matter of trust in the rewrite, we (i)~retain the verbatim original
reporter comment alongside the normalized question, so a reader can confirm that
no specific was invented or dropped, and (ii)~run a deterministic leak check that
flags any fix-type hard fact---a fixed version, fix pull request or commit, or
advisory URL---present in the question but absent from the pre-answer reporter
context: a fix the reporter stated earlier is legitimate context, whereas a fix
that could only have come from the answer is a leak. Flagged items are corrected
in the human pass.

\noindent\textbf{Thread-grounded answers.}
The gold answer is a faithful distillation of how the maintainer actually
answered, drawn strictly from the thread; the normalizer is forbidden to add
facts, fixes, or reasoning the thread does not contain. When the maintainer's
answer is terse or a pointer---``fixed in \#932, released 9.0.2''---the gold
answer reflects exactly that, labelled honestly as the maintainer's stated
resolution rather than inflated into a fuller explanation we cannot source. We
deliberately do \emph{not} fetch external advisories or patch diffs to enrich the
gold answer at this stage, as doing so would fabricate ground truth the thread
never contained. Instead, the artifacts the answer points to---advisory URLs and
the fix pull requests or commits---are recorded as \texttt{grounding\_sources};
at grading time the referenced pull-request and commit diffs are resolved and
supplied to the judge under the with-context conditions
(\S\ref{sec:experimental-setup}), so that a terse pointer-answer can still
corroborate a richer but correct model response without our having invented the
reference answer.

\noindent\textbf{Knowledge type: an explicit answerability axis.}
The most consequential annotation we attach records whether answering a pair
requires the repository at all. Each pair is labelled \emph{parametric}---a
security or engineering fact any expert could give from general knowledge or from
a public standard the question itself cites, which an LLM could answer with no
repository or document---or \emph{grounded}---resting on this project's specific
behaviour or releases, or on an external source the thread points to. This
distinction is what keeps the evaluation fair. The common pointer-answer ``fixed
in pull request \#932'' is unanswerable without the repository, so scoring a
no-context model on it would measure memorization of one project's pull-request
numbering rather than security capability. Recording knowledge type per pair
turns the no-context condition into a legitimate capability test on the
parametric pairs, makes the grounded pairs the locus of the artifact-context
analysis (\S\ref{sec:experimental-setup}), and turns high no-context accuracy on
grounded pairs into direct evidence of training-data contamination rather than
reasoning. The label is assigned by the normalizer from the answer's semantics and
then hardened by a deterministic rule: if the answer itself states a fix artifact
also recorded in \texttt{hard\_facts}---matched against the answer text, so that a
reporter-cited subject fact does not trigger it---the pair is grounded regardless
of the model's label, because such an artifact is project- or advisory-specific
information a model cannot produce unaided. The rule errs toward \emph{grounded}
by design: mislabelling a repository-internal answer as parametric would wrongly
admit an unanswerable question into the no-context capability test, whereas the
opposite error merely also routes a pair through the richer context conditions.

\noindent\textbf{Conservative decomposition.}
A single thread sometimes bundles distinct information needs---``is my usage
actually affected?'' and ``how do I remediate?''---each answered separately by the
maintainer. We split a thread into multiple items only when the reporter genuinely
raised distinct questions \emph{and} the maintainer's answer addresses more than
one of them; we never invent a sub-question the thread does not answer. This
restraint is the line between reframing a real implicit ask and fabricating
gradeable-looking questions the thread cannot ground---the latter being precisely
the synthetic-question construction our authenticity claim rejects.

\noindent\textbf{Outcome and verification.}
An initial pass over \dnum{50} verified threads yields \dnum{58} gradeable items,
with \dnum{5} threads decomposed into more than one; \dnum{53} items are grounded
and \dnum{5} parametric. The grounded-heavy split is expected: the
hard-verifiability filter (\S\ref{sec:verification}) selects for answers that rest
on a specific fix or advisory, leaving general-knowledge answers a minority. Every
item retains its verbatim original and is marked for human verification, and
\emph{two researchers independently review every normalized pair}. Reviewing the
full set rather than a sample lets us report inter-annotator agreement (Cohen's
$\kappa$) on the knowledge-type label and on the leak and context-completeness
judgments over the entire benchmark, and reconcile every disagreement before a
pair enters the evaluated set.

\subsection{Open Coding Toward a Question Taxonomy}
\label{sec:open-coding}

To characterize \emph{what} developers ask about, we apply inductive open coding
to the verified pairs. The coding unit is the verbatim Q\&A exchange. An LLM
coder (\texttt{gpt-5.4-mini}, accessed through a uniform LiteLLM interface) is
prompted to assign one or two short noun-phrase codes naming the security
concern at issue, given the issue title and body, the Q\&A summary, the full
thread, and the focal question and answer. As with detection, an LLM-assisted
first pass makes coding tractable and consistent at corpus scale while keeping a
human in control of the final code: the first author then reviewed every pair in
the same annotation interface, free to edit or replace any LLM-assigned code, so
the resulting codebook reflects human judgment rather than model output. The
first-author verification note is deliberately withheld from the coder to avoid
anchoring its codes to the acceptance rationale.

This pass produced \dnum{261} code instances over the \dnum{158} pairs
(\dnum{1--2} per pair), with \dnum{235} distinct codes---a fine-grained
vocabulary that reflects how specific developer security questions are
(e.g.\ ``missing subject alternative names,'' ``HTTPS redirect scheme
downgrade,'' ``transitive dependency vulnerability''). The single most frequent
code, \emph{transitive dependency vulnerability}, accounts for \dnum{12}
instances; the long tail of singleton codes confirms that the corpus is not
dominated by a handful of canned questions.

The next step is axial coding: collapsing the open codes into a smaller set of
higher-level categories that will anchor the per-category analysis in
RQ2--RQ3. Table~\ref{tab:axial} presents a \emph{preliminary, author-proposed}
grouping of the \dnum{235} codes into \dnum{eleven} candidate categories,
reported here to convey the shape of the emerging taxonomy. The two largest
groups---HTTP request and response security (SSRF, redirects, headers, cookies)
and TLS and certificate validation---together cover roughly half the corpus,
consistent with the HTTP-client and web-framework projects mined so far. We
stress that this grouping is not final: category names and boundaries will be
fixed through axial coding with an independent second coder, and we will report
inter-rater agreement (Cohen's $\kappa$) on a held-out sample, together with a
saturation analysis as the corpus grows.

\begin{table}[t]
\caption{Repositories in the current SecDevQA corpus. \textbf{Pairs} = Q\&A
  pairs retained after manual verification; \textbf{Hard} = subset whose answer
  carries at least one externally verifiable fact (CVE/GHSA ID, fixed version,
  fix commit/PR, advisory URL). Stars rounded to thousands (\dnum{2026-06-08}).}
\label{tab:repos}
\centering
\footnotesize
\setlength{\tabcolsep}{4pt}
\begin{tabular}{l l r r r}
\toprule
\textbf{Repository} & \textbf{Lang.} & \textbf{Stars} & \textbf{Pairs} & \textbf{Hard} \\
\midrule
psf/requests              & Python & \dnum{54.0k} & \dnum{24} & \dnum{14} \\
axios/axios               & JS     & \dnum{109k} & \dnum{23} & \dnum{19} \\
expressjs/express         & JS     & \dnum{69.1k} & \dnum{23} & \dnum{17} \\
pyca/cryptography         & Python & \dnum{7.6k} & \dnum{21} & \dnum{13} \\
urllib3/urllib3           & Python & \dnum{4.0k} & \dnum{14} & \dnum{6} \\
stripe/stripe-node        & JS     & \dnum{4.4k} & \dnum{13} & \dnum{5} \\
python-pillow/Pillow      & Python & \dnum{13.6k} & \dnum{12} & \dnum{12} \\
rails/rails               & Ruby   & \dnum{58.5k} & \dnum{11} & \dnum{9} \\
auth0/node-jsonwebtoken   & JS     & \dnum{18.2k} & \dnum{9} & \dnum{5} \\
fastapi/fastapi           & Python & \dnum{99.0k} & \dnum{8} & \dnum{5} \\
\midrule
\textbf{Total (\dnum{10} repos)} & & & \textbf{\dnum{158}} & \textbf{\dnum{105}} \\
\bottomrule
\end{tabular}
\end{table}

\begin{table}[t]
\caption{Artifact types the maintainer's answer drew on, across the \dnum{158}
  verified pairs (a pair may draw on several). Annotations are LLM-extracted and
  spot-checked; they drive the context conditions of the evaluation.}
\label{tab:artifacts}
\centering
\footnotesize
\setlength{\tabcolsep}{5pt}
\begin{tabular}{l r r}
\toprule
\textbf{Artifact type} & \textbf{Pairs} & \textbf{\%} \\
\midrule
documentation (docs, README, changelog) & \dnum{126} & \dnum{79.7} \\
code (repository source) & \dnum{92} & \dnum{58.2} \\
advisory (GHSA / OSV / project advisory) & \dnum{35} & \dnum{22.2} \\
pr\_data (PR diff or discussion) & \dnum{23} & \dnum{14.6} \\
issue\_tracker (linked issues) & \dnum{23} & \dnum{14.6} \\
dependency\_manifest & \dnum{19} & \dnum{12.0} \\
external\_reference (RFC, registry, etc.) & \dnum{14} & \dnum{8.9} \\
commit\_history (log, blame, commit) & \dnum{14} & \dnum{8.9} \\
cve\_cwe\_db (NVD/CWE/CVSS) & \dnum{7} & \dnum{4.4} \\
ci\_logs & \dnum{2} & \dnum{1.3} \\
\bottomrule
\end{tabular}
\end{table}

\begin{table*}[t]
\caption{\emph{Preliminary, author-proposed} axial grouping of the open
  codes into candidate higher-level categories. \textbf{Pairs} counts verified
  pairs touching the category (a two-code pair may fall in two categories, so
  the column sums to more than \dnum{158}). This grouping is a working draft to be
  finalized by axial coding with an independent second coder and inter-rater
  agreement; category names and boundaries are not yet fixed.}
\label{tab:axial}
\centering
\footnotesize
\begin{tabular}{l r p{0.62\linewidth}}
\toprule
\textbf{Candidate category} & \textbf{Pairs} & \textbf{Representative open codes} \\
\midrule
HTTP Request and Response Security & \dnum{40} & raw request body, credential leakage on redirect, CSRF token generation, double submit cookie, automatic redirect handling, absolute URL bypass \\
TLS and Certificate Validation & \dnum{38} & SSL certificate verification failure, OpenSSL version mismatch, ASN.1 parsing error, TLS SNI hostname handling, certificate verification failure, stricter SSL verification \\
Dependency and Supply-Chain Security & \dnum{26} & transitive dependency vulnerability, transitive dependency CVE, compromised package versions, dependency version update, prepare script install failure, false positive dependency vulnerability \\
Authentication and Session Integrity & \dnum{17} & webhook signature verification, unsigned JWT handling, question regarding jwt signing, signature verification bypass, token blacklist, JWT logout invalidation \\
Cryptographic Key and Algorithm Use & \dnum{15} & SM2 key algorithm mismatch, RSA key length requirement, cryptographic algorithm default mismatch, post-quantum support roadmap, encrypted RSA key loading, silent hash digest change \\
Input Handling and Denial of Service & \dnum{15} & memory exhaustion DoS, unsafe boundary generation, ReDoS vulnerability, ReDoS protection, regex route matching, array index limit \\
Vulnerability Report Triage & \dnum{14} & false positive report, scanner false positive, false positive vulnerability report, false vulnerability report, vulnerability database error, advisory database update \\
Patch, Release, and Compliance & \dnum{14} & fixed version announcement, SSRF vulnerability backport, safe version guidance, patched version marking, TLS deprecation support, SSRF fix availability \\
Secure Configuration and Defaults & \dnum{13} & deprecated URL parsing, opt-in security setting, missing config types, android network security config, initializer order misconfiguration, missing rate limiting \\
Platform and Version Compatibility & \dnum{10} & RAT on machine, issue tracker reference, local version parsing, Windows Python compatibility, permission denied error, missing attribute regression \\
Concurrency and Data Integrity & \dnum{7} & global state mutation, thread safety data race, shared block cache, deadlocked transaction state, writes outside transaction, thread safety issue \\
\bottomrule
\end{tabular}
\end{table*}


\section{Experimental Design}
\label{sec:experimental-setup}

\subsection{Research Questions}
\researchquestion{RQ1}{What categories of security concerns motivate developer
questions in open-source projects, and which artifact types do maintainers draw
on to answer them?}

\researchquestion{RQ2}{How accurately do frontier LLMs, open-weight LLMs, and
coding agents answer real developer security questions, broken down by question
category and knowledge type?}

\researchquestion{RQ3}{Which artifact types are necessary to answer which
categories of developer security questions? What is the marginal accuracy
contribution of each artifact type when access to it is selectively granted or
withheld?}

\researchquestion{RQ4}{When agents choose autonomously which artifacts to
consult, do they consult the artifact types maintainers actually needed---and
when they fail, is it because they never retrieved the right artifact or
because they reasoned incorrectly over it?}

RQ1 is answered by the benchmark construction itself (\S\ref{sec:open-coding});
the category-to-artifact map it produces is simultaneously an empirical finding
about developer security practice and the ground-truth prediction that RQ3
tests: if maintainers needed a given artifact type to answer questions of a
category, then granting a system access to that artifact type should improve
its accuracy on the category, and withholding it should specifically degrade
it.

\subsection{Evaluation at the Moment of Asking: Time-Capped Snapshots}
\label{sec:snapshot}

A naive with-context evaluation would hand the system the repository at its
current head---but for most items the fix commit, the fix pull-request
discussion, and often the advisory itself entered the repository \emph{after}
the question was asked. Against current state, ``answering'' degenerates into
retrieving the posted resolution. We instead freeze the world at the moment the
question was asked. For an item with report time $t_\theta$, the
\emph{time-capped snapshot} comprises: the repository working tree at the last
default-branch commit preceding $t_\theta$, with commit history truncated
there; the project's issues and pull requests created at or before $t_\theta$,
with later comments and reviews removed; and the GitHub security advisories
referencing the repository published at or before $t_\theta$. The source
thread itself is always excluded from the issue corpus---the system must never
read the conversation it is being tested on. No live network access is
provided under any condition:
a system that consulted today's NVD entry would observe the resolution and void
the cap.

Two properties of this construction deserve emphasis. First, resolution
identifiers that post-date $t_\theta$ are unavailable \emph{by construction},
which changes what can be graded (\S\ref{sec:grading}): for such items the task
is to produce the correct diagnosis, verdict, and remediation direction, not to
name a future release number. Second, when the resolution \emph{pre}-dates the
question---a duplicate report already answered, a fix already merged or
released, an advisory already published---it is legitimately discoverable
inside the snapshot, and locating it is precisely the task a maintainer performs
when triaging a new report. The benchmark therefore contains a natural
retrieval sub-task whose ground truth we did not have to synthesize.

\subsection{Systems Under Test}
\label{sec:systems}

We evaluate three system classes. Each agent backbone is also evaluated as a
bare LLM, so every with-context result has a paired no-context control from the
same model.

\textbf{Bare LLMs (no context).} The model receives only the synthesized
question $q$ and answers from parametric knowledge. We evaluate frontier
closed-weight models and open-weight models under identical prompts
(temperature~0 where the API permits).

\textbf{Reference agent (typed tools).} Our harness implements a tool-calling
agent whose tools are \emph{partitioned by artifact type}: file reading,
listing, and code search over the snapshot working tree (\textsf{code});
commit log and ancestor-guarded commit display (\textsf{commits}); keyword
search and full-text retrieval over the capped issue and pull-request corpora
(\textsf{issues}, \textsf{prs}); and search and retrieval over the capped
advisory set (\textsf{advisory}). Because each tool belongs to exactly one group, the set of
artifact types an agent consulted on an item is read directly off its
transcript as tool-call counts---the instrument RQ4 requires---and access to an
artifact type is granted or withheld by enabling or disabling its tool group,
the instrument RQ3 requires. The agent runs a bounded tool-calling loop
(at most \dnum{15} rounds) before committing to an answer; all transcripts are
released.

\textbf{Off-the-shelf coding agents.} To ground the agent results in systems
practitioners actually use, we additionally evaluate production coding agents
(Claude~Code and OpenCode) headlessly. Each item runs in a sandbox directory
materializing its snapshot as plain files: the repository tree exported at the
snapshot commit \emph{without} version-control metadata (so post-report history
is physically absent), the capped issue, pull-request, and advisory corpora as
structured data files, and a correspondingly truncated commit log. Because a
shell-equipped agent cannot be \emph{provably} prevented from reaching the
network, this condition is best-effort capped: the prompt forbids external
access, the sandbox contains nothing post-report, and we release full
transcripts; we treat the residual risk as a stated limitation
(\S\ref{sec:threats}) and rely on the typed-tool agent for the airtight
analysis.

\subsection{Context Conditions and Selective Artifact Provision}
\label{sec:conditions}

Let $\mathcal{G} = \{\textsf{code}, \textsf{commits}, \textsf{issues},
\textsf{prs}, \textsf{advisory}\}$ be the tool groups. The annotated artifact
types map onto these groups in the obvious way: code, dependency manifests,
and documentation are served by \textsf{code}; advisories and CVE/CWE database
entries by \textsf{advisory}; external references have no serving group, since
live external sources are excluded by the time cap. A \emph{condition} is a
subset $S \subseteq \mathcal{G}$ of groups the system may access:

\begin{itemize}
  \item $S = \emptyset$ --- \textbf{no context}: the bare-LLM condition;
  \item $S = \mathcal{G}$ --- \textbf{full snapshot}: the agent conditions of
        \S\ref{sec:systems}, run over the entire benchmark;
  \item $S = \mathcal{G} \setminus \{g\}$ (leave-one-out) and $S = \{g\}$
        (single-group) --- \textbf{selective provision}, run on a stratified
        subset of contextual items whose annotated artifact types span at least
        two groups, so that ablating one group leaves the others intact.
\end{itemize}

The selective-provision runs operationalize RQ3 as a controlled experiment:
necessity of group $g$ for category $k$ manifests as an accuracy drop under
$\mathcal{G} \setminus \{g\}$ relative to $\mathcal{G}$, and sufficiency as
accuracy under $\{g\}$ approaching that under $\mathcal{G}$
(\S\ref{sec:metrics}). The full-snapshot agent runs simultaneously supply the
RQ4 observational data: per item, the transcript yields the set of groups the
agent actually consulted, which we compare against the groups serving the
maintainer-annotated artifact types.

\subsection{Grading}
\label{sec:grading}

Model answers are free text, so grading combines a deterministic tier that is
exact where exactness is possible with an LLM-judge tier that is validated
where judgment is unavoidable. Both tiers are aware of the condition and the
knowledge type.

\textbf{Tier 1: condition-aware verifiable-fact matching.}
Each verifiable fact is matched in the response by a conservative, type-specific
identifier pattern (word-boundary--guarded version strings, PR-shaped
references, SHA prefixes of at least seven characters, exact CVE/GHSA/CWE
identifiers). Crucially, a fact is \emph{scored} only where it is knowable
under the condition: subject facts the reporter could cite are gradeable under
any condition, while a fix fact is gradeable only when the system had context
access \emph{and} the artifact behind the fact existed at report time (the fix
commit merged, the release tagged, the advisory published before $t_\theta$).
A no-context model is therefore never penalized for not knowing one project's
pull-request numbering, and no system is penalized for not naming an
identifier that did not yet exist---the failure mode that, left uncorrected,
turns a capability benchmark into a memorization probe. Unknowable facts are
recorded as \emph{not gradeable} and excluded from both numerator and
denominator.

\textbf{Tier 2: per-claim judge with a resolution oracle.}
An LLM judge---never the model under evaluation, with candidate identity
anonymized---scores the response against the item's prepared rubric
(\S\ref{sec:rubric-prep}), verdicting each criterion (a cause, a verdict, a
user-actionable remediation, an identifier) as met, partially met, or not met
and quoting the supporting response phrase so every verdict is spot-checkable. Because gold
answers are thread-grounded (\S\ref{sec:normalization}), a terse maintainer
pointer (``fixed in \#932, released 9.0.2'') yields a terse gold; to let a
richer-but-correct response earn credit such a gold cannot confirm, items
whose gold answer cites a fix pull request or commit are enriched at grading
time with the resolved change diff as a \emph{resolution oracle} for the judge.
Eligibility for this oracle is deliberately mechanical---an explicit fix PR or
commit reference, resolved by a single API call, with failures marked
unresolved and the item graded on the thread gold alone---and we report the
eligibility and resolution rates rather than chasing fixes through advisory
prose. The judge is condition-aware in one further respect: under agent
conditions, specific detail beyond the terse gold (affected ranges, root-cause
analysis, identifiers the agent verifiably retrieved) is expected enrichment,
and is flagged as hallucination only when it \emph{contradicts} the gold or the
question's premises; under no context, any concrete project-specific assertion
absent from the gold is flagged.

\textbf{Outcomes.} Each (item, system, condition) triple receives an outcome in
$\{\textsc{correct}, \textsc{partial}, \textsc{incorrect}\}$ by rolling the
per-criterion verdicts up along the rubric axes (\S\ref{sec:rubric-prep}):
correctness \emph{gates} the bucket---an item is \textsc{correct} only when every
correctness criterion is met, \textsc{incorrect} when any is missed or
contradicted, and \textsc{partial} otherwise (correctness only partially met, or
met with thin completeness coverage). Completeness criteria refine within the
bucket and are reported separately as a coverage rate, but never override
correctness, so we need no numeric criterion weights. Orthogonally, a boolean
\textsc{hallucinated} flags a response that asserts fabricated specifics. We keep hallucination separate rather than
folding it into \textsc{incorrect}: an answer can convey the right remediation
\emph{and} fabricate a CVE number, and the fabrication is the more
safety-relevant signal.

\textbf{Judge reliability.} We use an ensemble of $\ge 2$ judge models with
randomized presentation order; disagreements are flagged for human resolution.
Judge verdicts are cross-checked against the Tier-1 identifier matches, and
mismatches (e.g.\ a \textsc{correct} verdict with zero gradeable facts matched)
are audited. We calibrate against human labels on a stratified random sample of
\dnum{50--80} graded items and report Cohen's $\kappa$ (judge vs.\ human); the
judge prompt is revised and re-calibrated until $\kappa \ge 0.80$ before any
reported run.

\subsection{Metrics}
\label{sec:metrics}

Our primary metric is $\mathrm{Acc}(S, k)$, a system's accuracy under
condition $S$ on the items of question category $k$, alongside the analogous
partial-credit rate, the hallucination rate, and the Tier-1 verifiable-fact recall
over gradeable facts. All metrics are additionally split by knowledge type:
no-context capability claims are made on parametric items only, while contextual
items carry the artifact-context analysis. The marginal value of artifact
group $g$ for category $k$ is the leave-one-out contrast on the
selective-provision subset,
$\delta_g(k) = \mathrm{Acc}(\mathcal{G}, k) -
\mathrm{Acc}(\mathcal{G} \setminus \{g\}, k)$, and its sufficiency is
$\sigma_g(k) = \mathrm{Acc}(\{g\}, k) - \mathrm{Acc}(\emptyset, k)$, both over
contextual items. RQ1's category-to-artifact map predicts $\delta_g(k)$ to be
large exactly where maintainers needed an artifact type served by $g$; RQ3
tests that prediction. For RQ4 we report, per category, the agreement between
the groups the agent consulted and the groups serving the maintainer-annotated
artifact types, and decompose agent failures into \emph{retrieval failures}
(the agent answered incorrectly having never consulted some required group)
versus \emph{reasoning failures} (it consulted every required group and still
answered incorrectly).

\subsection{Negative Controls and Contamination Analysis}
\label{sec:contamination}

All threads are public and plausibly present in training corpora; a credible
security benchmark must measure, not assume away, memorization. We use three
instruments. First, the knowledge-type design itself: contextual items are
near-unanswerable without project access, so above-chance no-context accuracy
on contextual items is direct evidence of memorization, cleaner than any
post-hoc heuristic. Second, every item carries $t_\theta$, and results are
stratified by whether $t_\theta$ precedes each model's training cutoff;
markedly higher no-context accuracy on pre-cutoff items indicates recall rather
than reasoning. Third, for the typed-tool agent, the transcript makes
memorization visible at item granularity: a response naming a fix identifier
whose serving tool group was never called cannot have derived it from the
provided context. Approximately \dnum{10\%} of items are additionally tagged as
unanswerable from their annotated artifact types alone (requiring, e.g., an
external database lookup excluded by the time cap); these negative controls
measure confabulation separately from legitimate inference.

\section{Threats to Validity}
\label{sec:threats}

\noindent\textbf{Construct validity.}
Our evaluation items are \emph{synthesized} from raw threads---made
self-contained and stripped of their resolution---which could dilute the
authenticity claim or leak the answer. Both risks are addressed by construction
(\S\ref{sec:normalization}): the verbatim original reporter text is released
with every item; the deterministic leak check flags any fix identifier
present in a question that the reporter had not stated before the
answer, and we report the residual leak rate of the released set; and every
synthesized item is independently reviewed by two researchers, with Cohen's
$\kappa$ reported over the full benchmark for the knowledge-type, leak, and
context-completeness judgments and every disagreement reconciled before
evaluation. The knowledge-type label itself is a construct: mislabelling a
contextual pair as parametric would admit an unanswerable question into the
no-context capability test. The deterministic hardening rule
(\S\ref{sec:normalization}) biases errors in the safe direction, and
condition-aware gradeability ensures a labelling error cannot
convert an unknowable identifier into a scored miss. Finally, gold answers are
thread-grounded: a terse maintainer pointer yields a terse gold. We do not
inflate such golds; instead the grading-time resolution oracle
(\S\ref{sec:grading}) supplies the resolved fix diff to the judge, and we
report results separately for pointer-style and substantive-answer items so the
two populations are not conflated.

\noindent\textbf{Internal validity.}
Free-text grading uses an LLM judge, and verifiable facts ground but do not eliminate
judge error---a response may cite the correct CVE while contradicting the
verdict it implies. We mitigate with the controls of \S\ref{sec:grading}:
per-claim verdicts with quoted evidence rather than holistic scores, a
$\ge 2$-judge ensemble with flagged disagreements, regex cross-checks of
identifier verdicts, and human calibration to $\kappa \ge 0.80$ on a stratified
sample before any reported run. The resolution oracle is scoped by a mechanical
eligibility rule (an explicit fix PR or commit reference, one resolution
attempt); items outside it are graded on the thread gold alone, and we report
eligibility and resolution rates so the oracle's coverage is visible. Snapshot
assembly is itself an instrument that can err: the snapshot commit is resolved
on the default branch (a question may concern a release branch), and the issue and advisory
corpora are bounded by what our miners retrieved. We release snapshot
identifiers (commit SHAs, corpus counts, database snapshot dates) with every
run for re-verification.

\noindent\textbf{External validity.}
We mine only publicly visible artifacts. Security issues are often disclosed
privately and resolved under embargo, leaving just an advisory or changelog in
public; SecDevQA captures questions negotiated in the open and may
under-represent privately routed ones, so our claims are scoped to publicly
resolved security Q\&A. Repositories are selected for catalogued security
activity and popularity (\dnum{three} languages so far), which may bias toward
projects with formal advisory processes; we report per-repository results and
cap any repository's share of the corpus to keep the bias bounded and visible.
The hard-verifiability filter excludes question types that rarely yield
checkable identifiers (e.g.\ severity assessment, design rationale); deriving
the taxonomy from the full corpus, including non-hard pairs, keeps that gap
measurable rather than silent.

\noindent\textbf{Contamination and the time cap.}
Public threads may be memorized~\cite{stackrepoqa2025}, and the time cap
constrains a system's \emph{context}, not its \emph{weights}: a model whose
training data includes the resolved thread can reproduce the resolution under
any condition. We therefore measure memorization rather than assume its
absence, with three instruments (\S\ref{sec:contamination}): no-context
accuracy on contextual items, training-cutoff stratification over $t_\theta$, and
transcript-level detection for the typed-tool agent (a fix identifier asserted
without the serving tool group ever being called). A residual limitation
remains for off-the-shelf agents: a shell-equipped agent cannot be provably
prevented from reaching the network, so the external-agent condition is
best-effort capped---the sandbox physically contains no post-report data and
the prompt forbids external access---and we release its transcripts; airtight
claims rest on the typed-tool reference agent.

\section{Related Work}
\label{sec:related}

\noindent\textbf{Repository-level developer question answering.}
A recent line of benchmarks evaluates LLMs on questions about real codebases.
SWE-QA~\cite{sweqa2026} constructs repository-level questions from a taxonomy of
seed templates instantiated per repository with program analysis and an LLM,
pairing them with RAG-assisted, expert-revised answers and an LLM-as-judge
evaluation with human validation. Its questions are clean and controllable but
deliberately synthetic: no developer ever asked them. CoReQA~\cite{coreqa2025}
similarly targets repository-level comprehension with code-centric context.
StackRepoQA~\cite{stackrepoqa2025} moves toward authenticity by pairing Stack
Overflow questions with repositories, and reports a finding that motivates much
of our design: a substantial share of apparent model accuracy on public
developer threads reflects memorization of those threads rather than reasoning
over the repository. All three treat developer questions as a uniform
population, supply only code-internal context, and grade against answers whose
correctness is itself difficult to verify externally. SecDevQA differs on each
axis: its questions are verbatim-grounded in real maintainer exchanges, are
specifically security questions whose answers depend on artifacts \emph{outside}
the repository (advisories, vulnerability databases), carry externally
verifiable facts that anchor grading, and are evaluated under a time cap
that makes the memorization risk StackRepoQA observed measurable per item
(\S\ref{sec:contamination}) rather than merely acknowledged.

\noindent\textbf{Security benchmarks for LLMs and agents.}
Security-specific evaluations target tasks adjacent to ours but distinct from
it. SecVulEval~\cite{secvuleval2025} measures vulnerability \emph{detection} on
labelled code; CodeSecEval~\cite{codeseceval2024} measures secure code
\emph{generation}; SecureAgentBench~\cite{secureagentbench2025} evaluates
whether coding agents produce secure patches under repository tasks; and
SecQA~\cite{secqa2023} tests textbook security knowledge through
multiple-choice questions. These benchmarks share deterministic or
near-deterministic grading, but none evaluates the question-answering task
developers actually pose to maintainers---free-text, situation-specific, and
resolved by consulting heterogeneous artifacts---and none treats artifact
access as a controlled variable. SecDevQA is complementary: where
SecureAgentBench asks whether an agent can \emph{write} a secure change, we ask
whether it can \emph{answer} the security questions that precede and motivate
such changes, and which evidence it needs to do so.

\noindent\textbf{Developer security information behavior.}
That developers' security outcomes depend on the information sources they
consult is well established. Acar et al.~\cite{acar2016information} showed
experimentally that developers permitted to consult Stack Overflow produced
less secure code than those using official documentation, and that only a
minority of consulted threads contained secure advice; Fischer et
al.~\cite{fischer2017stackoverflow} traced insecure code snippets from such
threads into hundreds of thousands of deployed Android applications. This
literature measures the \emph{consequences} of answer quality but offers no
instrument for evaluating a new class of answerers. As LLMs and agents become
the first responders to developer security questions, SecDevQA provides that
instrument, and its category-to-artifact map (RQ1) extends this literature with
a quantitative account of \emph{which evidence} trustworthy answers rest on.

\noindent\textbf{Positioning.}
Table~\ref{tab:related} summarizes the comparison. The claim we make is
precise: no existing benchmark combines (a)~real security Q\&A mined from
developer forums, (b)~external advisory artifacts as a context variable,
(c)~fact-grounded, condition-aware grading, and (d)~systematic measurement of
which artifact types are necessary for which question categories. We do not
claim to be the first repository-level QA benchmark, nor the first security
evaluation of LLMs.

\begin{table}[t]
\caption{Comparison with related work.
  \textbf{Sec.}~=~security-specific corpus;
  \textbf{Adv.}~=~external advisory artifacts as context;
  \textbf{Grd.}~=~fact-grounded grading;
  \textbf{Abl.}~=~artifact-type ablation.
  \checkmark~full; P~partial; $\times$~none.}
\label{tab:related}
\centering
\footnotesize
\setlength{\tabcolsep}{4pt}
\begin{tabular}{l c c c c}
\toprule
\textbf{Paper} & \textbf{Sec.} & \textbf{Adv.} & \textbf{Grd.} & \textbf{Abl.} \\
\midrule
SWE-QA~\cite{sweqa2026}            & $\times$ & $\times$ & P          & $\times$ \\
StackRepoQA~\cite{stackrepoqa2025} & $\times$ & $\times$ & $\times$   & $\times$ \\
CoReQA~\cite{coreqa2025}           & $\times$ & $\times$ & P          & $\times$ \\
SecVulEval / CodeSecEval           & \checkmark & $\times$ & \checkmark & $\times$ \\
SecureAgentBench                   & \checkmark & $\times$ & \checkmark & $\times$ \\
SecQA                              & \checkmark & $\times$ & \checkmark & $\times$ \\
\midrule
\textbf{SecDevQA (ours)}           & \checkmark & \checkmark & \checkmark & \checkmark \\
\bottomrule
\end{tabular}
\end{table}


\bibliographystyle{IEEEtran}
\bibliography{references}

\end{document}